\documentclass[12pt]{article}
\usepackage[utf8]{inputenc}
\usepackage[T1]{fontenc}
\usepackage{amsmath,amssymb}
\usepackage{geometry}
\geometry{margin=2.5cm}

\begin{document}

\section*{Curs 3 -- Pagina 6}

\noindent\fbox{Teorema 1} \ Fie $(S,\varphi)$, $S\neq\varnothing$ şi $R$ o congruență pe $S$.

\bigskip
\noindent\underline{Atunci:} \ 1) $\varphi^* : S/R\times S/R \to S/R$ definită prin
\[
\varphi^*(\hat a,\hat b) = \varphi(a,b)
\]
s.n. legea de compoziție indusă pe $S/R$ de $\varphi$.

\begin{enumerate}
\item[(2)] $\varphi$ = asoc.\ (comut.) $\Rightarrow$ $\varphi^*$ = asoc.\ (comut.).
\item[(3)] $e\in S$ = elem.\ neutru pt.\ $\varphi$ $\Rightarrow$ $\hat e$ = el.\ neutru al operației $\varphi^*$.
\item[(4)] $a\in S$ este simetrizabil în rap.\ cu $\varphi$ $\Leftrightarrow$ $\hat a$ = simetrizabil în raport cu $\varphi^*$, şi $\hat a^{-1} = \widehat{a^{-1}}$.
\item[(5)] $(S,\varphi)$ = grup $\Rightarrow$ $(S/R,\varphi^*)$ = grup.
\end{enumerate}

\noindent $\big(\mathbb{Z}/n\mathbb{Z} = (\mathbb{Z}_n,+) = $ grup$\big)$

\bigskip
\noindent\underline{\textbf{Dem}}

\begin{enumerate}
\item[(1)] Nu se dem., e bine definit(ă).\footnote{ultimul cuvânt e neclar în original (arată ca ,,dif.'' în manuscris) -- probabil ,,definit(ă)''.}
\item[(2)] $\hat a\hat b = \widehat{ab} = \widehat{ba} = \hat b\hat a$.
\item[(3)] $\hat e\hat a = \widehat{ea} = \hat a = \widehat{ae} = \hat a\hat e$.
\item[(4)] $\hat a\hat a' = \widehat{aa'} = \hat e = \widehat{a'a} = \hat a'\hat a$.
\end{enumerate}

\[
p : S \to S/R
\]
\noindent $(S,\cdot)$ \hspace{2cm} $(S/R,\cdot)$

\[
p(ab) = p(a)\cdot p(b) \ \Leftrightarrow\ \widehat{ab} = \hat a\hat b.
\]

\[
\mathbb{Z} \to \mathbb{Z}_n
\]
\noindent $p(a) = \hat a$.

\[
p(a+b) = p(a)+p(b) \ \Leftrightarrow\ \widehat{a+b} = \hat a+\hat b.
\]
\[
p(ab) = p(a)\cdot p(b) \ \Leftrightarrow\ \widehat{ab} = \hat a\cdot\hat b.
\]

\end{document}
